\documentclass[12pt,a4paper]{article}
\usepackage[T1]{fontenc}
\usepackage[utf8]{inputenc}
\usepackage[spanish]{babel}
\usepackage{geometry}
\usepackage{booktabs}
\usepackage{xcolor}
\usepackage{listingsutf8}
\usepackage{hyperref}
\usepackage{enumitem}
\usepackage{titlesec}
\usepackage{parskip}
\usepackage{mdframed}
\usepackage{amssymb}

\geometry{margin=2.5cm}

\hypersetup{
  colorlinks=true,
  linkcolor=black,
  urlcolor=blue!70!black,
}

\lstdefinestyle{bash}{
  language=bash,
  basicstyle=\small\ttfamily,
  backgroundcolor=\color{gray!10},
  frame=single,
  rulecolor=\color{gray!40},
  breaklines=true,
  showstringspaces=false,
  inputencoding=utf8,
  extendedchars=true,
}

% Caja de advertencia
\newmdenv[
  linecolor=red!60!black,
  backgroundcolor=red!5,
  linewidth=1.5pt,
  topline=true, bottomline=true, leftline=true, rightline=true,
  innerleftmargin=8pt, innerrightmargin=8pt,
  innertopmargin=6pt, innerbottommargin=6pt,
]{advertencia}

% Caja informativa
\newmdenv[
  linecolor=blue!50!black,
  backgroundcolor=blue!5,
  linewidth=1pt,
  innerleftmargin=8pt, innerrightmargin=8pt,
  innertopmargin=6pt, innerbottommargin=6pt,
]{nota}

\title{%
  \textbf{Guía de reinstalación: CachyOS + Omarchy}\\[0.4em]
  \large Configuración post-instalación para TLOU Part~I y Bluetooth\\
  ASUS ROG Zephyrus G14 GA403UV
}
\author{pc-config}
\date{Mayo de 2026}

\begin{document}

\maketitle
\tableofcontents
\newpage

% ─────────────────────────────────────────────────────────────────────────────
\section{Contexto del hardware}

\begin{itemize}
  \item \textbf{Modelo}: ASUS ROG Zephyrus G14 GA403UV
  \item \textbf{iGPU}: AMD Radeon RX 780M (HawkPoint) --- driver \texttt{amdgpu},
        conectada directamente al panel interno.
  \item \textbf{dGPU}: NVIDIA GeForce RTX 4060 Laptop (8\,GB VRAM, Ada Lovelace
        AD107) --- driver \texttt{nvidia-open}, modo PRIME Offload.
  \item \textbf{Bluetooth/Wi-Fi}: MediaTek MT7922 --- chipset combinado,
        requiere parche DKMS para funcionar correctamente.
\end{itemize}

% ─────────────────────────────────────────────────────────────────────────────
\section{La regla más importante: versiones de NVIDIA}

\begin{advertencia}
\textbf{NUNCA actualizar los paquetes de userspace de NVIDIA más allá de 580.x.}
La versión 595 del Vulkan ICD (\texttt{nvidia-utils}) provoca artefactos de
renderizado severos en \emph{The Last of Us Part~I}: bloques en rejilla sobre
los modelos de personajes, humo roto y vegetación corrupta. Este comportamiento
es reproducible e inaceptable.
\end{advertencia}

La combinación correcta y verificada es:

\begin{center}
\begin{tabular}{lll}
\toprule
\textbf{Paquete} & \textbf{Versión} & \textbf{Motivo} \\
\midrule
\texttt{nvidia-open-dkms} & 595.x & Corrige el bug CrashCAT (GSP queue overflow) \\
\texttt{nvidia-utils} & 580.x & El ICD 595 causa artefactos visuales \\
\texttt{lib32-nvidia-utils} & 580.x & Ídem \\
\texttt{nvidia-settings} & 580.x & Ídem \\
\texttt{opencl-nvidia} & 580.x & Ídem \\
\texttt{lib32-opencl-nvidia} & 580.x & Ídem \\
\bottomrule
\end{tabular}
\end{center}

El módulo del kernel al 595 y el userspace al 580 generan mensajes de
``API mismatch'' en \texttt{dmesg}, pero Vulkan funciona correctamente para
gaming. Este split está verificado con una sesión de más de 20 minutos en TLOU.

\begin{nota}
\textbf{¿Por qué el kernel necesita 595?}
El driver 580 tiene un bug en el firmware GSP que provoca un desbordamiento
en \texttt{crashcat\_queue\_v1.c:66} bajo carga de compilación de shaders.
Esto causa un cuelgue o crash del juego a los pocos minutos. La versión 595
del módulo del kernel corrige este bug.
\end{nota}

% ─────────────────────────────────────────────────────────────────────────────
\section{Pasos antes de instalar Omarchy}

Realizar estos pasos con el sistema base de CachyOS ya instalado y con
acceso a internet, \textbf{antes} de ejecutar el instalador de Omarchy.

% ──────────────────────────────────────────
\subsection{Paso 1: Bloquear el userspace de NVIDIA en pacman}

Editar \texttt{/etc/pacman.conf} e insertar o completar la línea
\texttt{IgnorePkg} \textbf{antes de instalar ningún paquete de NVIDIA}:

\begin{lstlisting}[style=bash]
sudo nano /etc/pacman.conf
\end{lstlisting}

Buscar la línea \texttt{IgnorePkg} (puede estar comentada con \texttt{\#}).
Descomentarla y dejarla así:

\begin{lstlisting}[style=bash]
IgnorePkg = nvidia-utils lib32-nvidia-utils nvidia-settings \
            opencl-nvidia lib32-opencl-nvidia
\end{lstlisting}

\textbf{No} incluir \texttt{nvidia-open-dkms} en esta lista.

% ──────────────────────────────────────────
\subsection{Paso 2: Instalar los drivers de NVIDIA con chwd}

\begin{lstlisting}[style=bash]
sudo chwd -a
\end{lstlisting}

Esto instalará \texttt{nvidia-open-dkms} y \texttt{nvidia-utils} a la versión
más reciente disponible en los repos de CachyOS (actualmente 595.x para ambos).

\begin{nota}
Si \texttt{chwd} ya fue ejecutado durante la instalación de CachyOS y los
drivers ya están presentes, omitir este paso y continuar con el Paso~3.
\end{nota}

% ──────────────────────────────────────────
\subsection{Paso 3: Degradar el userspace a 580}

Instalar la herramienta \texttt{downgrade}:

\begin{lstlisting}[style=bash]
sudo pacman -S downgrade
\end{lstlisting}

Degradar cada paquete del userspace. El comando mostrará una lista de
versiones disponibles; seleccionar \textbf{580.95.05-1} en cada caso:

\begin{lstlisting}[style=bash]
sudo downgrade nvidia-utils
sudo downgrade lib32-nvidia-utils
sudo downgrade nvidia-settings
sudo downgrade opencl-nvidia
sudo downgrade lib32-opencl-nvidia
\end{lstlisting}

Verificar que el resultado sea correcto:

\begin{lstlisting}[style=bash]
pacman -Q nvidia-open-dkms nvidia-utils lib32-nvidia-utils \
           nvidia-settings opencl-nvidia lib32-opencl-nvidia
\end{lstlisting}

Salida esperada (puede haber variación menor en el número de compilación):

\begin{lstlisting}
nvidia-open-dkms      595.71.05-2
nvidia-utils          580.95.05-1
lib32-nvidia-utils    580.95.05-1
nvidia-settings       580.95.05-1
opencl-nvidia         580.95.05-1
lib32-opencl-nvidia   580.95.05-1
\end{lstlisting}

% ──────────────────────────────────────────
\subsection{Paso 4: Restaurar el firmware GSP 595}

Al degradar \texttt{nvidia-utils} de 595 a 580, pacman elimina los archivos
de firmware \texttt{nvidia/595.71.05/gsp\_ga10x.bin} y
\texttt{nvidia/595.71.05/gsp\_tu10x.bin}. El módulo del kernel 595 los necesita
para inicializar la GPU. Hay que extraerlos del paquete que quedó en la caché.

\begin{lstlisting}[style=bash]
# Verificar que el paquete de nvidia-utils 595 este en cache
ls /var/cache/pacman/pkg/nvidia-utils-595*

# Extraer el firmware a un directorio temporal
TMPDIR=$(mktemp -d)
bsdtar -xf /var/cache/pacman/pkg/nvidia-utils-595.71.05-*-x86_64.pkg.tar.zst \
       -C "$TMPDIR" usr/lib/firmware/nvidia/595.71.05/

# Copiar al destino definitivo
sudo mkdir -p /usr/lib/firmware/nvidia/
sudo cp -r "$TMPDIR/usr/lib/firmware/nvidia/595.71.05" \
           /usr/lib/firmware/nvidia/
rm -rf "$TMPDIR"

# Verificar
ls -lh /usr/lib/firmware/nvidia/595.71.05/
\end{lstlisting}

La salida debe mostrar dos archivos:
\begin{lstlisting}
gsp_ga10x.bin   ~73 MB
gsp_tu10x.bin   ~30 MB
\end{lstlisting}

\begin{nota}
Si el paquete \textbf{no está en caché} (por ejemplo, en una instalación
limpia donde \texttt{chwd} descargó 595 pero la caché fue limpiada),
descargarlo sin instalarlo:
\begin{lstlisting}[style=bash]
sudo pacman -Sw --nodeps nvidia-utils   # descarga sin instalar
\end{lstlisting}
Luego repetir el bloque de extracción anterior.
\end{nota}

% ──────────────────────────────────────────
\subsection{Paso 5: Reconstruir la imagen UKI / initramfs}

Con el firmware ya en disco, regenerar las imágenes del kernel para que
el firmware quede embebido en el UKI y esté disponible en el arranque
temprano (antes de montar el sistema de ficheros raíz):

\begin{lstlisting}[style=bash]
sudo mkinitcpio -P
\end{lstlisting}

Cuando termine (tarda 2--5 minutos), verificar que el firmware está en el UKI:

\begin{lstlisting}[style=bash]
sudo lsinitcpio /boot/EFI/Linux/omarchy_linux.efi 2>/dev/null \
     | grep "nvidia/595"
\end{lstlisting}

Salida esperada:
\begin{lstlisting}
usr/lib/firmware/nvidia/595.71.05/
usr/lib/firmware/nvidia/595.71.05/gsp_ga10x.bin
usr/lib/firmware/nvidia/595.71.05/gsp_tu10x.bin
\end{lstlisting}

Si los archivos aparecen, el UKI está correcto.

\begin{nota}
En instalaciones con Limine y múltiples kernels (como CachyOS), el comando
\texttt{mkinitcpio -P} genera UKIs para todos los kernels instalados y
preguntará si ejecutar \texttt{limine-mkinitcpio}. Responder \texttt{Y}.
El UKI activo por defecto es \texttt{omarchy\_linux.efi}.
\end{nota}

% ──────────────────────────────────────────
\subsection{Paso 6: Instalar el parche DKMS para Bluetooth (MT7922)}

El driver \texttt{btmtk} del kernel devuelve \texttt{-EINVAL} cuando el
firmware del MT7922 responde con 7 bytes en lugar de los 9 esperados al
comando FUNC\_CTRL enable. El módulo DKMS \texttt{btmtk-fix} corrige esto
usando un poll de estado como fallback.

El código fuente del módulo está en el repositorio de configuración de este
equipo. Desde el directorio del repositorio:

\begin{lstlisting}[style=bash]
# El modulo fuente debe estar en /usr/src/btmtk-fix-1.0/
# (copiarlo alli si no esta)
sudo dkms add btmtk-fix/1.0
sudo dkms install btmtk-fix/1.0 --force

# Verificar
dkms status | grep btmtk
\end{lstlisting}

Salida esperada: \texttt{btmtk-fix/1.0, <kernel-version>: installed}

% ──────────────────────────────────────────
\subsection{Paso 7: Reiniciar y verificar antes de Omarchy}

\begin{lstlisting}[style=bash]
sudo reboot
\end{lstlisting}

Tras el reinicio, verificar desde la TTY o terminal:

\begin{lstlisting}[style=bash]
# La GPU debe haberse inicializado (debe existir /dev/nvidia0)
ls /dev/nvidia*
# Esperado: /dev/nvidia0  /dev/nvidiactl  /dev/nvidia-uvm  ...

# El modulo del kernel debe ser 595
cat /proc/driver/nvidia/version

# Bluetooth debe funcionar
bluetoothctl show
\end{lstlisting}

Si \texttt{/dev/nvidia0} no aparece, revisar \texttt{dmesg} buscando errores
de firmware:

\begin{lstlisting}[style=bash]
sudo dmesg | grep -i "nvidia\|gsp\|firmware" | grep -i "fail\|error"
\end{lstlisting}

Si aparece \texttt{gsp\_ga10x.bin failed with error -2}, el firmware no
está en el UKI --- repetir el Paso~5.

% ─────────────────────────────────────────────────────────────────────────────
\section{Instalar Omarchy}

Una vez verificado que NVIDIA y Bluetooth funcionan:

\begin{lstlisting}[style=bash]
bash <(curl -fsSL https://omarchy.dev/install)
\end{lstlisting}

Seguir las instrucciones del instalador. Omarchy configura Hyprland, Waybar
y las variables de entorno necesarias para NVIDIA en Wayland
(\texttt{NVD\_BACKEND=direct}, \texttt{LIBVA\_DRIVER\_NAME=nvidia}, etc.).

% ─────────────────────────────────────────────────────────────────────────────
\section{Configuración de Steam y TLOU Part~I}

\subsection{Instalar Steam}

\begin{lstlisting}[style=bash]
sudo pacman -S steam
\end{lstlisting}

O usar el script del repositorio de configuración si está disponible.

\subsection{Instalar Proton GE}

Desde Steam: \textbf{Ajustes → Compatibilidad → Activar Steam Play para
todos los títulos} y seleccionar Proton GE. Alternativamente con ProtonUp-Qt:

\begin{lstlisting}[style=bash]
sudo pacman -S protonup-qt   # o flatpak install ProtonUp-Qt
\end{lstlisting}

Abrir ProtonUp-Qt, seleccionar ``Steam'' y descargar la versión más reciente
de \textbf{GE-Proton}.

\subsection{Instalar VKD3D-Proton 3.0.0}

\begin{advertencia}
\textbf{Versión específica requerida.} TLOU Part~I requiere VKD3D-Proton
exactamente en la versión \textbf{3.0.0}. Versiones posteriores pueden
presentar regresiones con este juego.
\end{advertencia}

Descargar \texttt{vkd3d-proton-3.0.0.tar.zst} desde:
\texttt{https://github.com/HansKristian-Work/vkd3d-proton/releases/tag/v3.0.0}

Extraer en el directorio de Proton GE activo:

\begin{lstlisting}[style=bash]
# Sustituir GE-ProtonXX-YY por la version instalada
PROTON_DIR="$HOME/.steam/root/compatibilitytools.d/GE-ProtonXX-YY"
tar -xf vkd3d-proton-3.0.0.tar.zst -C "$PROTON_DIR/files/lib64/wine/" \
    --strip-components=1
\end{lstlisting}

\begin{nota}
ProtonUp-Qt también puede instalar versiones específicas de VKD3D-Proton
directamente desde su interfaz, lo que es más cómodo.
\end{nota}

\subsection{Opciones de lanzamiento en Steam para TLOU Part~I}

En Steam: clic derecho en el juego → \textbf{Propiedades → Opciones de
lanzamiento}. Escribir únicamente:

\begin{lstlisting}
%command%
\end{lstlisting}

No son necesarias variables especiales (\texttt{DXVK\_FILTER\_DEVICE\_NAME},
\texttt{PROTON\_NO\_ESYNC}, etc.). El juego usa Vulkan directamente a través
del Vulkan ICD de NVIDIA.

% ─────────────────────────────────────────────────────────────────────────────
\section{Mantenimiento tras actualizaciones del kernel}

Cuando se actualiza el kernel (p.ej., de \texttt{6.18.x} a \texttt{6.19.x}),
los módulos DKMS deben reconstruirse para el nuevo kernel:

\begin{lstlisting}[style=bash]
sudo dkms autoinstall
\end{lstlisting}

Esto reconstruye tanto \texttt{nvidia-open-dkms} como \texttt{btmtk-fix} para
el nuevo kernel.

Después, regenerar el UKI:

\begin{lstlisting}[style=bash]
sudo mkinitcpio -P
\end{lstlisting}

Verificar que el firmware siga en el UKI (ver Paso~5) y reiniciar.

\begin{nota}
El directorio \texttt{/usr/lib/firmware/nvidia/595.71.05/} \textbf{no está
gestionado por pacman} (fue copiado manualmente). Las actualizaciones del
sistema no lo tocan. Solo necesita reponerse si se reinstala el sistema o si
se actualiza \texttt{nvidia-open-dkms} a una versión que use un directorio de
firmware diferente (p.ej., 596.x.x usaría \texttt{nvidia/596.x.x/}).
\end{nota}

Si se instala un nuevo kernel mayor que requiera una versión de firmware
distinta, repetir el Paso~4 (extracción del firmware) con la versión nueva
del paquete correspondiente.

% ─────────────────────────────────────────────────────────────────────────────
\section{Referencia rápida: IgnorePkg definitivo}

El bloque completo de \texttt{IgnorePkg} que debe estar en
\texttt{/etc/pacman.conf} de forma permanente:

\begin{lstlisting}[style=bash]
IgnorePkg = nvidia-open-dkms nvidia-utils lib32-nvidia-utils \
            nvidia-settings opencl-nvidia lib32-opencl-nvidia
\end{lstlisting}

\begin{advertencia}
\texttt{nvidia-open-dkms} también está en \texttt{IgnorePkg} para evitar
actualizaciones automáticas no supervisadas. Si se desea actualizar el módulo
del kernel en el futuro, quitarlo temporalmente de \texttt{IgnorePkg},
actualizar, extraer el nuevo firmware, reconstruir el UKI y volver a añadirlo.
\end{advertencia}

\end{document}
